% Reviewed FF.30 copy, integrated in F_section.tex after independent review. Do not input this duplicate.
\begin{claim}{FF.30}{Proved}\label{clm:FFpadicsubfields}
Let \(F\subset\overline{\mathbb Q}_p\) be a field
of finite transcendence degree over \(\mathbb Q\).
Suppose that, for the induced valuation normalized
by \(v(p)=1\),
\[
                  v(F^\times)\subset
                         p^{-a}\mathbb Z_{(p)}
                    \quad\text{for some }a\geq0.       \tag{FF.30.1}
\]
Then \(\mathcal R(F)\) holds. In particular this
applies when the induced valuation on \(F\) is
discrete. No henselianity assumption is required.
\end{claim}
\begin{proof}
Choose a finitely generated field
\(D\subset F\) containing a transcendence basis
of \(F/\mathbb Q\). Its closure inside
\(\overline{\mathbb Q}_p\) is a finite extension
\(M/\mathbb Q_p\): it is generated topologically
by \(\mathbb Q_p\) and the finitely many
\(\mathbb Q_p\)-algebraic generators of \(D\).
Let \(K_D\) be the chosen strict henselization
field of the valuation ring of \(D\), inside
its algebraic closure in
\(\overline{\mathbb Q}_p\), and let
\(T_D\) be the maximal tame extension of \(K_D\)
there. Completion and the chosen prolongation
identify
\[
               G_{K_D}=I_M,\qquad G_{T_D}=P_M.         \tag{FF.30.2}
\]
Indeed the henselization of \(D\) is dense in
\(M\) and has the same finite separable extensions
by the Krasner--Hensel comparison proved in
Claim~\ref{clm:FFpdetector}. Strict henselization
then fixes inertia, and adjoining the prime-to-\(p\)
roots of a uniformizer fixes wild inertia.

We verify that Claims~\ref{clm:FFpdetector},
\ref{clm:FFstrictreconstruct} and
\ref{clm:FFtamereconstruct} apply to finite
extensions of \(K_D\) and \(T_D\), respectively.
The first proof uses only a henselian discrete
valuation of characteristic zero, algebraically
closed residue field of characteristic \(p\),
its completion, semistable finite-part
\(p\)-divisible groups, and Tate's theorem over
the completed DVR. These hypotheses hold here.
The exact multiplicative kernel follows from
the same discrete binomial estimate, with
\(\mu_{p'}\) supplied by strict henselization.

In the arithmetic-equivariance step of
Claim~\ref{clm:FFstrictreconstruct}, replace
its number fields \(D_i\) by finite extensions
of the fixed \(D\) containing the finitely
many model coefficients, detector coordinates,
and actual point values. Such fields suffice
because all the fields under consideration are
algebraic over \(D\). Their completions
\(M_i\) are finite \(p\)-adic fields, and
(FF.30.2) gives exactly the required normal
inertia subgroups. The finite-point conjugation
argument is therefore unchanged. The Picard
1-motive step uses compactness only for
\(\operatorname{Pic}^0(M_i)\), not for the
uncompleted model field. Its resulting torsion
or zero equality for a \(D_i\)-rational Picard
point descends from \(M_i\) to \(D_i\), since
field extension is faithfully flat. A principal
boundary divisor is represented by a normalized
function over \(D_i\) by the same Hilbert~90
argument. Thus the unit power, resultant,
zero-functional argument, scalar removal and
finite-extension descent all apply.

For \(T_D\), the norm proof of the multiplicative
kernel and finite-image descent of the functional
space use only the pro-prime-to-\(p\) quotient
of inertia by wild inertia. The valuation and
tame-root calculations in
Claim~\ref{clm:FFtamereconstruct} likewise use
only the DVR \(K_D\). The argument gives
\(\mathcal R(K)\) for every finite extension
\(K/T_D\).

Finally \(F/D\) is algebraic. The bounded
\(p\)-denominator argument of
Claim~\ref{clm:FFboundedwild}, applied with
\(K_D\) in place of \(K_p\), shows that
\(FT_D/T_D\) is finite. The normalization of
its value group may include the fixed
ramification index of \(M/\mathbb Q_p\);
this changes the finite bound by a fixed factor
and does not affect the argument. Therefore
\(\mathcal R(FT_D)\) holds. The extension
\(FT_D/F\) is algebraic because \(D\subset F\)
and \(T_D/D\) is algebraic. Algebraic descent
proves \(\mathcal R(F)\).
\end{proof}
